% !TeX root = ../main.tex
\section{Evaluation Methodology}
\label{sec:evaluation}

\subsection{Simulation}

\textbf{Simulation.}
We extend MGPUSim~\cite{sun2019mgpusim} with distributed page placement, a
packet-switched mesh, RDMA endpoints, detailed cache and HBM timing, and the
instrumentation required to model the complete memory path of a wafer-scale
GPU.  We model 48 GPMs connected by a $7\times7$ mesh, with the remaining tile
hosting the central I/O and control components.  Each link provides 768~GB/s
of bandwidth and incurs 32 cycles of switch latency, consistent with prior
multi-chip and wafer-scale studies
\cite{arunkumar2017mcm,feng2024barre}.  Each GPM represents one quarter of an
AMD MI100-class GPU.  Four GPMs collectively provide the compute, cache,
memory capacity, and memory bandwidth of one MI100
\cite{techpowerup_mi100}.  This scaling keeps cycle-level simulation of the
48 GPM system tractable while preserving its aggregate resource ratios.

\textbf{Cache and memory hierarchy.}
Each GPM contains 32 CUs organized as eight shader arrays.  Every CU has a
private 16~KB, four-way L1 vector cache with a 10-cycle access, 16 MSHRs, and
at most 160 concurrent transactions.  Each shader array also has a 16~KB
scalar cache and a 32~KB instruction cache.  The GPM-wide 4~MB L2 is divided
into four independently accessed 1~MB partitions.  Each partition is 16-way
associative, accepts 16 requests per cycle, contains 64 MSHRs, and charges a
10-cycle directory lookup to both
hits and misses.  Four memory controllers connect the L2 to 8~GB of local HBM
with 1.23~TB/s aggregate bandwidth.  Each RDMA endpoint processes eight
transfers per cycle, charges 10 cycles for an endpoint traversal, and tracks
up to 64 requester and 64 owner transactions.  Table~\ref{table:BaselineConfig}
summarizes the modeled system.

\textbf{CuPath modeling.}
The MCM-GPU design introduces a GPM-side, remote-only L1.5 cache between L1
and L2 and explores reallocating L2 capacity to that new level
\cite{arunkumar2017mcm}.  CuPath does not add an L1.5 cache, a victim cache,
or another independently addressable cache.  Baseline and CuPath therefore
use identical L1 and L2 capacities, associativities, ports, MSHRs, and lookup
latencies.  CuPath adds one metadata-only Cuckoo Filter and a bounded 16-line
M1 response buffer per L2 partition, together with bounded local and remote
prediction state.  The response buffer holds an adjacent line returned by a
paired HBM access until a matching demand consumes it or replacement removes
it.  A retained remote response is instead inserted as a clean line through
the existing L2 fill and replacement path.  Exact cache tags, MSHRs, and RDMA
transaction tables remain authoritative.  All
configurations also use the same TLBs and 4~KB physical page mapping, so every
reported difference comes from cache, HBM, RDMA, and network processing below
address translation.

\textbf{Execution scope.}
Each experiment continues until 78{,}600 workgroups have completed across the
48-GPM system.  This represents approximately 51 completed workgroups for each
of the wafer's 1,536 CUs and is sufficient to capture steady execution beyond
the initial transient \cite{baruah2020griffin}.  We terminate at a workgroup
boundary because prior studies show that workgroup completion rates stabilize
after warmup \cite{avalos2021principal,liu2023photon}.  The threshold does not
change kernel launches, workgroup IDs, partitioning, or GPU assignment.  A
workload with fewer workgroups runs to completion, and all configurations use
the same termination rule.

\begin{table}[t]
  \centering
  \caption{Baseline wafer-scale GPU and CuPath configuration.}
  \label{table:BaselineConfig}
  \footnotesize
  \begin{tabularx}{\columnwidth}{@{}l>{\raggedright\arraybackslash}X@{}}
    \toprule
    \textbf{Module} & \textbf{Configuration} \\
    \midrule
    Wafer & 48 GPMs, 1.0~GHz, $7\times7$ mesh \\
    Compute & 32 CUs/GPM, eight shader arrays/GPM \\
    L1 vector & 16~KB/CU, 4-way, 10 cycles,
                16 MSHRs \\
    L1 scalar & 16~KB/shader array, 4-way, 16 MSHRs \\
    L1 instruction & 32~KB/shader array, 4-way, 16 MSHRs \\
    Shared L2 & 4~MB/GPM, four partitions, 16-way, 64 MSHRs/partition,
                16 req./cycle/partition, 10-cycle lookup \\
    HBM & 8~GB/GPM, four controllers, 1.23~TB/s, BL4 \\
    RDMA endpoint & 8 transfers/cycle, 10-cycle processing,
                    64 outstanding/direction \\
    Mesh link & 768~GB/s, 32-cycle switch latency, 16-B flit \\
    \midrule
    \multicolumn{2}{l}{\textit{CuPath additions}} \\
    Cache capacity & No L1.5 or new addressable cache; baseline L1/L2 unchanged \\
    Typed Filter & One shared array/L2 partition, 4K buckets, four slots/bucket,
                   13-bit fingerprint \\
    Filter ports & 1-cycle lookup/update, 16 lookups and 16 updates/cycle/partition \\
    M1 state & 64 observed regions and 16 returned lines/L2 partition \\
    Remote predictor & 256 entries/requesting GPM \\
    M1 & Filter-guided paired HBM access \\
    M2 & Up to eight lines/packet, 64 packet descriptors/requesting GPM \\
    M3 & No separate history; best-effort clean fill into existing L2 \\
    \bottomrule
  \end{tabularx}
\end{table}

\subsection{Ablation Configurations}

We run five configurations from the same frozen simulator binary.
\emph{Baseline} disables every CuPath transformation.  \emph{M1} enables the
complete local request-pairing path: reliable \textsc{Resident} negatives
bypass guaranteed-miss L2 tag lookups, and typed membership state admits at
most one predicted sibling into the triggering miss's paired HBM access.
\emph{M2} uses \textsc{Remote-Pending} to guide exact same-line
deduplication and \textsc{Remote-Batch} to locate joinable, work-conserving
home-GPM/page packets. Predicted remote candidates can use only a free position
in a Filter-identified and exactly confirmed existing batch;
requesting-GPM L2 reuse is off.  \emph{M3} enables \textsc{Seen}/
\textsc{Resident}-guided requesting-GPM L2 probe gating and reuse without remote
batching or candidate generation.  \emph{Complete} enables M1--M3.  Row-aware DRAM
reordering remains disabled in all five configurations.

In the complete configuration, the stages operate on the surviving request
population.  A requesting-GPM L2 termination never enters requesting-GPM RDMA, and a
coalesced RDMA waiter does not create a network or home-GPM memory request.  We therefore place
counters at each boundary and report both eliminated work and end-to-end
performance, rather than treating the three stage speedups as additive.

\subsection{Hardware Overhead}

We estimate CuPath's storage and area overhead with CACTI~7.0
\cite{balasubramonian2017cacti} at 32~nm.  Each L2 partition has a 4K-bucket
Typed Filter with four slots per bucket.  Each slot contains a 13-bit
fingerprint, three type bits, one valid bit, and a four-bit reference count,
requiring 42~KiB per Filter.  We model each Filter as 16 banks with two read
ports and one write port per bank.  In the conflict-free case, this organization
supports 16 two-bucket lookups and 16 updates per cycle.  The 256-entry remote
predictor is conservatively budgeted at 64~B per entry, or 16~KiB per GPM.
The four Filters and predictor therefore require 184~KiB per GPM, equal to
4.49\% of the 4-MiB L2 data capacity.

CACTI estimates 0.358~mm$^2$ for each Filter and 0.036~mm$^2$ for the remote
predictor.  The resulting CuPath area is 1.469~mm$^2$ per GPM.  A complete
4-MiB, 16-way L2, modeled as four 1-MiB partitions and including data arrays,
tag arrays, and peripheral circuitry, occupies 54.510~mm$^2$.  CuPath therefore
adds 2.70\% area relative to the modeled L2.  Filter and predictor accesses take
0.523~ns and 0.241~ns, respectively, and fit within the 1-ns cycle.  The estimate
includes the arrays and their modeled ports, but excludes hashing, fingerprint
comparators, arbitration, bank conflicts, and CuPath control logic.

\subsection{Workloads and Execution}

We evaluate CuPath with 14 benchmarks from Hetero-Mark
\cite{sun2016hetero}, AMD APP SDK~\cite{amd_opencl_guide},
SHOC~\cite{danalis2010scalable}, and DNNMark~\cite{dong2017dnnmark}.
The suite has also been used by prior multi-GPU studies
\cite{wang2025oasis,li2023idyll,li2021improving,feng2024barre} and spans
local and remote accesses with both streaming and repeated data use.
Table~\ref{table:Benchmark} reports these two workload properties, the
workgroups actually admitted and retired in the formal run and the workload
allocation footprint measured immediately after allocation from the
simulator's 4-KB page counters.  We derive the labels from the latest Baseline
run.  \textbf{Remote} and \textbf{Local} identify whether most memory requests
are served by a remote or local home GPM.  \textbf{Repeated} indicates that L1
and L2 cache hits together serve more than half of read accesses, whereas
\textbf{Streaming} indicates that most reads pass through both cache levels
and continue toward memory.  Baseline and Complete allocate the same number of
pages and complete the same number of workgroups for every workload.

We choose the largest hundreds-of-MB or GB-scale input that remains tractable
in cycle-level simulation, following prior multi-GPU and wafer-scale studies
\cite{baruah2020griffin,li2023trans,xu2025exploring}.  Every footprint exceeds
the capacity of a single GPM's private and shared caches, forcing requests to
exercise local HBM and, under distributed page placement, remote home-GPM paths.

The runner records every launch, partition, observed workgroup count, and a
hash of the global workgroup set.  We reject a comparison if these identities
differ across configurations.

\begin{table}[t]
  \centering
  \caption{Benchmarks, access location, data use, completed workgroups, and footprint.}
  \label{table:Benchmark}
  \scriptsize
  \setlength{\tabcolsep}{3pt}
  \begin{tabularx}{\columnwidth}{lccXr}
    \toprule
    \textbf{Benchmark} & \textbf{Access} & \textbf{Data use} &
    \textbf{Completed WGs} & \textbf{Footprint} \\
    \midrule
    AES  & Local  & Repeated  & 78,600 & 1,024~MiB \\
    BT   & Local  & Streaming & 78,600 & 1,024~MiB \\
    FWT  & Local  & Streaming & 78,600 & 1,024~MiB \\
    FFT  & Local  & Streaming & 78,600 & 1,024~MiB \\
    FIR  & Local  & Repeated  & 78,600 & 1,024~MiB \\
    FWS  & Local  & Repeated  & 78,600 & 1,024~MiB \\
    I2C  & Remote & Repeated  & 78,600 & 1,029~MiB \\
    KM   & Local  & Repeated  & 73,728 & 390~MiB \\
    MM   & Remote & Repeated  & 78,600 & 1,024~MiB \\
    MT   & Remote & Streaming & 32,400 & 1,013~MiB \\
    PR   & Remote & Repeated  & 78,600 & 1,024~MiB \\
    RELU & Local  & Streaming & 78,600 & 2,560~MiB \\
    SC   & Local  & Repeated  & 78,600 & 512~MiB \\
    SPMV & Local  & Streaming & 78,600 & 1,081~MiB \\
    \bottomrule
  \end{tabularx}
\end{table}
